\documentclass[11pt]{article}

\usepackage[margin=1in]{geometry}
\usepackage{amsmath,amssymb,amsthm}
\usepackage{mathtools}
\usepackage{xcolor}
\usepackage{listings}
\usepackage{hyperref}
\hypersetup{colorlinks=true,linkcolor=blue!50!black,urlcolor=blue!50!black,citecolor=blue!50!black}

% ---- Lean code listing style -------------------------------------------------
\definecolor{leankw}{rgb}{0.50,0.10,0.45}
\definecolor{leancm}{rgb}{0.30,0.45,0.30}
\lstdefinestyle{lean}{
  basicstyle=\ttfamily\small,
  keywordstyle=\color{leankw}\bfseries,
  commentstyle=\color{leancm}\itshape,
  columns=fullflexible,
  keepspaces=true,
  breaklines=true,
  showstringspaces=false,
  morekeywords={theorem,lemma,def,noncomputable,by,rw,simp,simp_rw,have,exact,intro,
    refine,calc,apply,set,congr,ring,field_simp,linarith,nlinarith,gcongr,fun,let,
    Measure,map,open,namespace,instance},
  morecomment=[l]{--},
  morecomment=[s]{/-}{-/},
}
\lstset{style=lean}

\newtheorem{theorem}{Theorem}
\newtheorem{lemma}{Lemma}
\theoremstyle{remark}
\newtheorem*{remark}{Remark}

\newcommand{\R}{\mathbb{R}}
\newcommand{\E}{\mathbb{E}}
\newcommand{\Prob}{\mathbb{P}}
\newcommand{\eps}{\varepsilon}
\newcommand{\norm}[1]{\lVert #1 \rVert}
\DeclareMathOperator{\mgf}{mgf}

\title{\textbf{A Machine-Checked Proof of the\\ Johnson--Lindenstrauss Lemma in Lean~4}}
\author{Formalized over mathlib}
\date{}

\begin{document}
\maketitle

\begin{abstract}
We present a complete formalization, in the Lean~4 proof assistant on top of
\texttt{mathlib}, of the Johnson--Lindenstrauss (JL) lemma: a random Gaussian projection
of $n$ points from $\R^d$ into $\R^m$ with $m = O(\log n / \eps^2)$ dimensions preserves
all pairwise squared distances to within a factor $(1\pm\eps)$, with high probability and
\emph{independently of the ambient dimension $d$}. The development compiles end to end
with no \texttt{sorry} and depends only on the three standard foundational axioms. We
describe the mathematical argument, its decomposition into seven independently verified
modules, and the principal Lean techniques used at each step.
\end{abstract}

\section{What we prove, and the plan}

\paragraph{The statement.} Fix $0<\eps<1$ and $n$ distinct points $x_1,\dots,x_n\in\R^d$.
Let $A$ be an $m\times d$ random matrix with i.i.d.\ $N(0,1/m)$ entries. The
Johnson--Lindenstrauss lemma asserts that, provided $m \gtrsim \log n / \eps^2$,
\[
  (1-\eps)\,\norm{x_i-x_j}^2 \;\le\; \norm{A x_i - A x_j}^2 \;\le\; (1+\eps)\,\norm{x_i-x_j}^2
  \qquad\text{for all } i\ne j,
\]
simultaneously, except with probability at most $\delta$. Concretely, the failure
probability is at most $\delta$ as soon as $2n^2\exp(-m\eps^2/12)\le\delta$, i.e.
$m\ge 12\,(\log(2n^2)+\log(1/\delta))/\eps^2$.

\paragraph{The strategy.} We follow the elementary proof of Dasgupta and
Gupta~\cite{dasgupta}. Because $A$ is Gaussian, for a \emph{fixed} vector $v$ the image
$Av$ is itself Gaussian, $Av\sim N(0,(\norm{v}^2/m)I_m)$, so $\norm{Av}^2$ is a scaled
chi-squared variable. Concentration of the chi-squared distribution (via its moment
generating function and the Chernoff bound) controls the distortion of a single vector;
a union bound over the $\binom n2$ difference vectors $x_i-x_j$ finishes the proof. The
argument splits into seven steps, realized as seven Lean modules
(Table~\ref{tab:modules}).

\begin{table}[h]
\centering
\begin{tabular}{ll}
\hline
\textbf{Module} & \textbf{Result}\\
\hline
\texttt{LogBounds}   & \S\ref{sec:log} elementary $\log$ inequalities\\
\texttt{GaussianMGF} & \S\ref{sec:gmgf} $\E[e^{tZ^2}]=(1-2t)^{-1/2}$, $Z\sim N(0,1)$\\
\texttt{Basic}, \texttt{ChiSq} & \S\ref{sec:chisq} the law $\chi^2_m$ and its MGF $(1-2t)^{-m/2}$\\
\texttt{TailBounds}  & \S\ref{sec:tails} Chernoff tail bounds for $\chi^2_m$\\
\texttt{Projection}  & \S\ref{sec:proj} single-vector concentration\\
\texttt{UnionBound}  & \S\ref{sec:union} union bound $\Rightarrow$ JL\\
\hline
\end{tabular}
\caption{The seven modules and their results.}
\label{tab:modules}
\end{table}

\paragraph{The formal theorem.} The random projection is modelled by its $m$ independent
standard-Gaussian \emph{rows} $G=(G_1,\dots,G_m)$, $G_k\in\R^d$; the projected squared
norm of $v$ is $(1/m)\sum_k\langle v,G_k\rangle^2=\norm{Av}^2$ for $A=G/\sqrt m$. The law
of $G$ is \texttt{rowLaw m d}, and $\mu.\mathtt{real}\,S=(\mu S).\mathtt{toReal}$.

\begin{lstlisting}
theorem johnsonLindenstrauss (n d : N) (m : N) (m_pos : 0 < m)
    (x : Fin n -> EuclideanSpace R (Fin d)) (hx : Function.Injective x)
    (eps delta : R) (he_pos : 0 < eps) (he_lt : eps < 1) (hd_pos : 0 < delta)
    (hsuff : 2 * (n:R)^2 * rexp (-(m:R) * eps^2 / 12) <= delta) :
    (rowLaw m d).real
      {G | exists i j, i != j /\
        ((1+eps) * ||x i - x j||^2 <= (1/(m:R)) * sum k, <x i - x j, G k>^2
         \/ (1/(m:R)) * sum k, <x i - x j, G k>^2 <= (1-eps) * ||x i - x j||^2)}
      <= delta
\end{lstlisting}

\section{Logarithm bounds}\label{sec:log}

The Chernoff exponents are $\eps-\log(1+\eps)$ and $-(\eps+\log(1-\eps))$. Integrating the
pointwise inequality $\tfrac1{1+x}\le 1-x+x^2$ over $[0,\eps]$,
\[
  \eps-\log(1+\eps) \;=\; \int_0^\eps\!\Big(1-\tfrac1{1+x}\Big)\,dx
  \;\ge\; \int_0^\eps (x-x^2)\,dx \;=\; \tfrac{\eps^2}2-\tfrac{\eps^3}3,
\]
and symmetrically $\eps+\log(1-\eps)\le-\eps^2/2$. For $0<\eps<1$ both simplify to the
clean form used downstream:
\[
  \eps-\log(1+\eps)\ \ge\ \tfrac{\eps^2}{6},\qquad
  \eps+\log(1-\eps)\ \le\ -\tfrac{\eps^2}{6}.
\]

\begin{remark}
The frequently quoted bound $\eps-\log(1+\eps)\ge\eps^2/3$ is \emph{false} near $\eps=1$
(at $\eps=1$ the left side is $1-\log2\approx0.307<1/3$). The sharp constant of this form
valid on all of $(0,1)$ is $\eps^2/6$, since
$\tfrac{\eps^2}2-\tfrac{\eps^3}3\ge\tfrac{\eps^2}6\iff 2\eps^2(1-\eps)\ge0$. We use
$\eps^2/6$ throughout; this is one place where formalization corrected a folklore
constant.
\end{remark}

In Lean the comparison is an interval-integral inequality. The antiderivative
$\int_0^\eps\tfrac1{1+x}=\log(1+\eps)$ comes from a substitution
(\texttt{integral\_comp\_add\_left}) and \texttt{integral\_one\_div}; the comparison is
\texttt{intervalIntegral.integral\_mono\_on}, with integrability discharged by
\texttt{Continuous.intervalIntegrable} and the pointwise step by \texttt{nlinarith}.

\section{The MGF of a squared standard normal}\label{sec:gmgf}

For $Z\sim N(0,1)$ and $t<1/2$, completing the square inside the Gaussian integral,
\[
  \E[e^{tZ^2}] = (2\pi)^{-1/2}\!\int e^{tx^2}e^{-x^2/2}\,dx
  = (2\pi)^{-1/2}\!\int e^{-(\frac12-t)x^2}\,dx
  = (2\pi)^{-1/2}\sqrt{\tfrac{\pi}{\frac12-t}} = (1-2t)^{-1/2}.
\]
The Lean proof passes to the density via \texttt{integral\_gaussianReal\_eq\_integral\_smul},
applies mathlib's Gaussian integral \texttt{integral\_gaussian}
($\int e^{-bx^2}=\sqrt{\pi/b}$), and discharges the radical identity by normalizing both
sides to $\sqrt{(1-2t)^{-1}}$ (\texttt{Real.sqrt\_inv}, \texttt{Real.sqrt\_mul},
\texttt{field\_simp}, \texttt{Real.sqrt\_eq\_rpow}).

\section{The chi-squared distribution and its MGF}\label{sec:chisq}

The chi-squared law with $m$ degrees of freedom is the pushforward of the standard
Gaussian on $\R^m$ under $x\mapsto\norm x^2$:
\begin{lstlisting}
noncomputable def chiSq (m : N) : Measure R :=
  (stdGaussian (EuclideanSpace R (Fin m))).map (fun x => ||x||^2)
\end{lstlisting}
Writing the standard Gaussian on $\R^m$ as the pushforward of $\bigotimes N(0,1)$ makes
$\norm g^2=\sum_i g_i^2$, so the MGF \emph{factorizes}:
\[
  \E[e^{t\norm g^2}] = \prod_{i=1}^m \E[e^{tg_i^2}] = \big((1-2t)^{-1/2}\big)^m
  = (1-2t)^{-m/2}.
\]
The key mathlib ingredient is \texttt{integral\_fintype\_prod\_eq\_prod}, the
product-measure factorization $\int\prod_i f_i(x_i)\,d(\bigotimes\mu)=\prod_i\int f_i$;
the $m$-fold product is collapsed with \texttt{rpow\_natCast} and \texttt{rpow\_mul}.

\begin{lemma}[\texttt{chiSq\_mgf}]
For $t<1/2$, $\ \mgf\,\mathrm{id}\,(\chi^2_m)\,t = (1-2t)^{-m/2}$.
\end{lemma}

\section{Chernoff tail bounds}\label{sec:tails}

Applying the Chernoff method to \texttt{chiSq\_mgf} and optimizing $t$
($t=\eps/(2(1+\eps))$ for the upper tail, $t=-\eps/(2(1-\eps))$ for the lower),
\[
  \Prob(\chi^2_m\ge(1+\eps)m)\le e^{-\frac m2(\eps-\log(1+\eps))}\le e^{-m\eps^2/12},
  \quad
  \Prob(\chi^2_m\le(1-\eps)m)\le e^{\frac m2(\eps+\log(1-\eps))}\le e^{-m\eps^2/12}.
\]
The Lean tail bounds invoke \texttt{measure\_ge\_le\_exp\_mul\_mgf} and its lower
counterpart; integrability of $e^{tX}$ is free, since the MGF is finite and nonzero so the
integrand cannot be non-integrable (\texttt{integral\_undef}, contrapositive). The
exponent identity reduces, after \texttt{Real.rpow\_def\_of\_pos}, to
\texttt{field\_simp; ring}, and the simplification to $e^{-m\eps^2/12}$ chains through
\S\ref{sec:log}.

\section{Single-vector concentration}\label{sec:proj}

For fixed $v\ne0$, $\norm{Av}^2\sim(\norm v^2/m)\chi^2_m$. We model this law as
\texttt{projNormSq m v} $=(\chi^2_m).\mathtt{map}(z\mapsto(\norm v^2/m)z)$ and reduce each
distortion event to a chi-squared tail by change of variables
(\texttt{Measure.map\_apply}), obtaining
\[
  \Prob\big(\norm{Av}^2\notin[(1-\eps)\norm v^2,(1+\eps)\norm v^2]\big)\le 2e^{-m\eps^2/12}.
\]
The two-sided event is split with \texttt{measureReal\_union\_le} and the two
\S\ref{sec:tails} bounds applied via \texttt{gcongr}.

\section{The union bound and the JL lemma}\label{sec:union}

Two ingredients complete the proof.

\paragraph{The marginal law (the technical crux).} With $m$ i.i.d.\ standard-Gaussian
rows, the projected norm of a fixed $v$ has law exactly \texttt{projNormSq m v}:
\begin{lstlisting}
lemma rowLaw_map_projSq (m d : N) (v : EuclideanSpace R (Fin d)) :
    (rowLaw m d).map (fun G => (1/(m:R)) * sum i, <v, G i>^2) = projNormSq m v
\end{lstlisting}
This is established by a chain of pushforwards: each coordinate
$\langle v,G_k\rangle\sim N(0,\norm v^2)$ via
\texttt{IsGaussian.map\_eq\_gaussianReal} (mean $0$, variance $\norm v^2$ from
\texttt{variance\_dual\_stdGaussian} and \texttt{innerSL\_apply\_norm}); these lift
coordinatewise through \texttt{Measure.pi\_map\_pi}; a rescaling
(\texttt{gaussianReal\_map\_const\_mul}) and the identity $\sum(\norm v\,y_i)^2=\norm
v^2\sum y_i^2$ collapse the result onto \texttt{chiSq}, reusing
\texttt{map\_pi\_eq\_stdGaussian}.

\paragraph{The union bound.} The bad event is contained in the union, over the $\le n^2$
distinct pairs, of the per-pair distortion events. By sub-additivity
(\texttt{measureReal\_biUnion\_finset\_le}) and the per-pair bound (\S\ref{sec:proj} via
the marginal),
\[
  \Prob(\text{bad}) \le \sum_{i\ne j} 2e^{-m\eps^2/12} \le n^2\cdot2e^{-m\eps^2/12}
  = 2n^2 e^{-m\eps^2/12} \le \delta.
\]
Injectivity of the points gives $x_i-x_j\ne0$ (\texttt{sub\_ne\_zero}), as the per-pair
bound requires. This yields the theorem of \S1.

\section{Assessment}

\paragraph{Trust.} \texttt{\#print axioms johnsonLindenstrauss} reports exactly
\texttt{[propext, Classical.choice, Quot.sound]} --- the standard \texttt{mathlib}
foundation, with no \texttt{sorryAx}. Every intermediate lemma, including the marginal,
was checked the same way.

\paragraph{The bound.} The explicit sufficient condition $2n^2e^{-m\eps^2/12}\le\delta$
gives $m=O(\log n/\eps^2)$, independent of $d$ --- the entire point of the lemma. The
constant $12$ is honest, tracing through the $\eps^2/6$ logarithm bound and the factor
$1/2$ in the chi-squared exponent.

\paragraph{Engineering.} The proof is seven small files ($\approx500$ lines), each
compiled independently and mirroring one step of the argument. Small, named lemmas were
preferred to monolithic tactic blocks at every turn: they are easier to state, to prove,
to read, and to repair.

\begin{thebibliography}{9}
\bibitem{dasgupta} S.~Dasgupta and A.~Gupta. \emph{An elementary proof of a theorem of
Johnson and Lindenstrauss}. Random Structures \& Algorithms 22(1):60--65, 2003.
\bibitem{jl} W.~B.~Johnson and J.~Lindenstrauss. \emph{Extensions of Lipschitz mappings
into a Hilbert space}. Contemp.\ Math.\ 26:189--206, 1984.
\bibitem{mathlib} The mathlib Community. \emph{The Lean mathematical library}. CPP 2020.
\end{thebibliography}

\end{document}
